%titolo{I canederli}
%autore{}
%album{}
%tonalita{}
%famiglia{Altre}
%gruppo{}
%momenti{}
%identificatore{i_canederli}
%data_revisione{2026_10_28}
%trascrittore{Francesco Endrici}
%video{}
\beginsong{I canederli}
\begiverse
Per fare dei canéderli
col bròdo e col ragù
se ciapa del prezémolo
e se lo taia su
farina òvi e zìgole
lugàneghe col spéck
pan vècio sénza mìgole
e n tòch de formài sgnec.
Se fà balòtole col pangratà
tre quatro frégole de ai pestà
na mèza chìcchera de lat e vin
èco i canéderli de noi trentìn.
\endverse
\beginverse
Enséma coi canéderli
noialtri ghe magnan
na tésa de lugàneghe
e 'n tòch de smacafàm
ne pias polénta e fìnferli
la mòsa con el lat
capussi con le scódeghe
che vanza fór dal piàt.
Per la maiòlica inoltre ghé
capussi e cràuti su per Piné
g'avén le tròte zò lì a Toblìn
col vino santo e el marzemìn.
\endverse
\beginverse
Gh'è po' la Val dei Mocheni
da Perzen nando 'n su.
se 'n contrerà Sant'Òrsola, 
Fieròzz e pò Palù;
gh'aven anche le zigole 
e qualche patatàr,
d'invèrno le luganeghe 
e schnaps da trinkenàr.
\endverse
\beginverse
Andando su a Sant’Orsola
dai reumi se guariss
coi bagni “terapeutici”
(almeni cossita i diss…);
nei boschi gh’è le giasene
per chi le vol binar,
che spezialmentee i stitici
le pol miracolar.
\endverse
\endsong